Versicherungs und Finanzmarktprodukte


Lebens

Nennen sie wesentliche Produktformen der privaten Lebens und Rentenversicherung

INSERT IMAGE 17 DAV GUIDE
Kapitalversicherung
(einmalige Leistungszahlung)

Kapitalbildende
Lebensversicherung

(Reine) Erlebensfall-Vers.



Systematisieren und erläutern Sie den technischen Leistungsumfang in
der privaten Lebens- und Rentenversicherung.

Der technische Leistungsumfang der Lebensversicherungsprodukte umfasst drei
Kategorien:
Im Zentrum stehen zunächst die sogenannten Versicherungsgarantien, die aus
der Übertragung biometrischer Risiken vom Versicherten auf den Versicherer ent-
stehen. Hierzu zählen insb. Todesfallschutz, Ablaufleistung im Erlebensfall.

Typischerweise enthalten Lebens- und Rentenversicherungsprodukte auch Fi-
nanzgarantien. Im klassischen Geschäftsmodell werden die Finanzgarantien
durch die Verwaltung des Deckungsstocks (Sicherungsvermögen) sichergestellt.
Fondsgebundene Versicherungsverträgen können auf Finanzgarantien völlig
verzichten, das Anlagerisiko trägt dann ausschließlich der Versicherungsnehmer.

Bei hybriden Lebensversicherungsprodukten werden
die Finanzgarantien teilweise im Deckungsstock und am Kapitalmarkt erzeugt. Die
klassische Lebensversicherung umfasst eine Garantieverzinsung pro Jahr (sog.
Höchstrechnungszins) sowie eine dem Grunde aber nicht der Höhe nach garan-
tierte Überschussbeteiligung (laufende und Schlussüberschussbeteiligung.)

Nicht zuletzt aufgrund der hohen Kosten einer jährlichen Zinsgarantie werden zunehmend Lebensversiche-
rungsprodukte angeboten, die eine endfällige Kapitalgarantie enthalten
Wobei das Prinzip gilt, je höher das Garantieniveau, desto geringer sind die
Wertentwicklungschancen.


Welche Bedeutung hat das einzelvertragliche versicherungstechnische
Äquivalenzprinzip für die prospektive Kalkulation des Deckungskapi-
tals?


Die Kalkulation der Lebensversicherungsprämien folgt dem versicherungstech-
nischen Äquivalenzprinzip.
Rechnungsgrundlagen in der Lebensversicherung sind neben dem Kalkulationszinssatz 
Sterbe- bzw. Überlebenswahrscheinlichkeiten, Stornowahrscheinlichkeiten und Zuschläge für
bestimmte Kostenbestandteile.

Die vorsichtige Wahl der Rechnungsgrundlagen
erster Ordnung (implizite Sicherheitszuschläge) führt mit hoher Wahrscheinlich-
keit zur Überschussentstehung beim Versicherer.


Eine am natürlichen Lebensprozess anknüpfende Prämienzahlung hätte für den
Versicherungsnehmer zur Folge, dass er zur Absicherung des Todesfallrisikos zu
Beginn der Vertragslaufzeit eine relativ geringe Periodenprämie zu entrichten
hätte, die dann im Zeitablauf sukzessive ansteigt (natürliche Prämie).

Im Grundmodell der Lebensversicherung entrichtet der Versicherungsnehmer hinge-
gen eine im Zeitablauf konstante Periodenprämie (tatsächliche Prämie).

Während dieser Phase bildet der Versicherer in Höhe des übersteigenden Anteils eine Kapitalre-
serve (Deckungsrückstellung).

Die Deckungsrückstellung entspricht in beiden
Fällen der positiven Differenz zwischen dem Barwert der erwarteten künftigen Leis-
tungszahlungen des Versicherers und dem Barwert der erwarteten künftigen Prä-
mienzahlungen des Versicherungsnehmers (prospektive Ermittlungsmethode).




Nennen und erläutern Sie wesentliche Produktformen der privaten
Krankenversicherung.


Gegenstand der Krankenhaustagegeldversicherung ist die
Zahlung des vereinbarten Krankenhaustagegeldes im Falle der medizinisch not-
wendigen stationären Heilbehandlung. Die Krankentagegeldversicherung er-
setzt den als Folge von Krankheit oder Unfall durch Arbeitsunfähigkeit verursach-
ten Verdienstausfall durch das vereinbarte Krankentagegeld. Im Fall der Pflegebe-
dürftigkeit deckt die Pflegekrankenversicherung Aufwendungen im vereinbar-
ten Umfang für die Pflege der versicherten Person (Pflegekostenversicherung) das
vereinbarte Tagegeld zu (Pflegetagegeldversicherung).

Auch in Deutschland können Mitglieder der gesetzlichen Krankenversicherung (GKV) private Zu-
satzversicherungen erwerben. Für nicht sozialversicherungspflichtige Personen
besteht in Deutschland die europaweit einzigartige Möglichkeit, substitutiven privaten 
Krankenversicherungsschutz zu erwerben.

Die substitutive private Krankheitskostenvollversicherung dient der Absicherung sämtlicher Kosten für 
ambulante und stationäre Heilbehandlung. Darüber hinaus können auch zahnärztliche Behandlungskosten eingeschlossen werden.

(ADD SOME POINTS FROM SLIDE 34 FROM KPMG)

Rechnungsgrundlagen sind in der Krankenversicherungsauf-
sichtsverordnung (KVAV) festgelegt und umfassen im Wesentlichen die sog. Kopf-
schäden, den Rechnungszins, die Ausscheideordnung (Sterblichkeit und Storno),
den Sicherheitszuschlag, weitere Zuschläge zur Deckung von Abschluss-, Verwal-
tungs- und Regulierungskosten sowie Zuschläge für die erfolgsunabhängige RfB.


Was verstehen wir unter einem „Kopfschadenprofil“ und welche Bedeu-
tung hat dessen Veränderung auf die Prämienkalkulation der PKV?

Die in der Kalkulation zugrunde gelegten periodischen Krankheitskosten (Kopfschäden) 
steigen mit zunehmendem Alter der versicherten Person an. Eine an dem
natürlichen Krankheitskostenverlauf angepasste Prämienzahlung würde dement-
sprechend ebenfalls periodisch ansteigen (natürliche Prämie). Dem Kalkulati-
onsprinzip der Lebensversicherung folgend wird jedoch eine periodisch konstante
Prämienzahlung vereinbart (tatsächliche Prämie). Dies führt in analoger Weise
zur Kapitalbildung (Alterungsrückstellung) in den ersten Versicherungsjahren
und zu einer entsprechenden Auflösung der Reserven gegen Ende der Vertrags-
laufzeit.


Erkennt der Versicherer im Zeitablauf, dass die der Kalkulation zugrunde liegenden
erwarteten Krankheitskosten (kalkulierte Kopfschäden) nicht ausreichen,
muss er die periodisch zu leistenden Beiträge an den tatsächlichen Risikoverlauf anzupassen.
Die einem individuellen Vertrag zuzurechnende Kostensteigerung ist unabhängig vom zwi-
schenzeitlich erreichten Gesundheitsstand des Versicherten, eine Gesundheitsprü-
fung erfolgt nicht. In zahlreichen Fällen wirken sich Kostensteigerungen in Abhän-
gigkeit der erwarteten Restlebensdauer eines Versicherungsnehmers überpropor-
tional auf die Höhe der Beitragsanpassung aus.


Erläutern Sie versicherungszweigübergreifend wesentliche Merkmale
der Schaden-/ Unfall-versicherung (Produktcharakteristika).


Der Begriff Kompositversicherung ist, mit Ausnahme der Krankenversicherung,
die Sammelbezeichnung für alle Versicherungsarten der Schaden- und Unfallver-
sicherung zusammen. Schadenversicherungen decken das konkrete Interesse
des Versicherten bzw. des geschädigten Dritten, mithin dem finanziellen Ersatz
des Sachschadens (Sachversicherung) sowie der Gefahrenabwehr von Haftungsri-
siken (Haftpflichtversicherungen).
Die private Unfallversicherung zählt hingegen zur Personenversicherung, 
gedeckt ist das abstrakte Interesse des Versicherungsnehmers, 
das durch vereinbarte Versicherungssumme bestimmt ist.

Das technische Risiko in der Sachversicherung ist durch die vertragliche Fixierung
der versicherten Sache (z.B. Wohngebäude, Kraftfahrzeug etc.), der ver-
sicherten Gefahren (z.B. Feuer) sowie der versicherten Schäden (Sachschaden)
aber auch indirekten Schäden, wie z.B. Aufräumkosten.
Insb. in den Zweigen der Schadenversicherung werden Verträge typischerweise
mit einer Laufzeit von einem Jahr abgeschlossen.



Nennen und charakterisieren Sie die bedeutendsten Versicherungs-
zweige in der Kompositversicherung.

Gemessen am Bruttoprämienvolumen ist die Kraftfahrtversicherung (Kfz-Haft-
pflicht, Kfz-Vollkasko- und -teilkaskoversicherung etc.)
Mit deutlichem Abstand folgen die Haftpflichtversicherung (Allgemeine Haft-
pflichtversicherung (AHV), Berufshaftpflichtversicherungen)
die private Unfallversicherung, die verbundene Wohngebäu-
deversicherung.


Charakterisieren Sie kapitalbildende Altersvorsorgeprodukte hinsicht-
lich ihrer Vermögensanlage, des Vermögensaufbaus und des Vermö-
gensverzehrs.


Mit Ausnahme des selbstgenutzten Wohneigentums erfolgt die
Vermögensanlage in allen beschriebenen Fällen kollektiv. Ursächlich hierfür ist die
sog. Losgrößen- und Fristentransformation, wonach zahlreiche kleinere Bei-
träge in einem Topf gebündelt und dann in größere Anlageobjekte kollektiv inves-
tiert werden.

Im Versicherungsbereich erfolgt die Vermögensbildung hingegen
durch die kollektive Verwaltung des Deckungsstocks an dessen Wertentwicklung
die Versicherungsnehmer unabhängig von ihrem Eintrittszeitpunkt (Vertragsab-
schluss) in identischer Weise beteiligt werden. Allein im Fall der privaten Renten-
versicherung ist auch der Vermögensverzehr an den kollektiv verwalteten De-
ckungsstock geknüpft

INSERT TABLE FROM SLIDE 26



Erläutern Sie die Idee und Wesensmerkmale des Bausparens.

Kern der Idee des Bausparens ist die Organisation eines Kollektivs von Sparern
und Darlehensnehmern, die mit dem Ziel der Immobilienfinanzierung in einen gemeinsamen Kapitalstock sparen 
und aus diesem zu einem späteren Zeitpunkt ein zinsgünstiges Baudarlehen beanspruchen. 
Im Kollektiv der Bausparer unterscheidet man daher zwischen Sparern und Darlehensnehmern.

In den Kapitalstock fließen entsprechend Sparbeiträge sowie Zins- und Tilgungszahlungen der Darlehens-
nehmer, darüber hinaus auch Erträge aus der Investition des Teils des kollektiven
Kapitalstocks, der noch nicht als Baudarlehen vergeben wurde.

Im Bausparvertrag werden zwischen Bausparer und Bausparkasse dem jeweili-
gen Bauspartarif entsprechend die Bausparsumme, die Ansparkonditionen (z.B.
Guthabenzins, Regelsparbeitrag, Mindestzeiten), die Zuteilungsvoraussetzungen
sowie die Darlehenskonditionen (Darlehenszins, Laufzeiten, Tilgungsbeiträge, Si-
cherheiten etc.) fest vereinbart.

Zunächst verpflichtet sich der Bausparer (Anleger) in der Ansparphase, den vertraglich vereinbarten Teil der Bausparsumme einzuzahlen. 
Das Bausparguthaben setzt sich zusammen aus der Summe der geleisteten Einzahlungen, ggf. ergänzt um staatliche Förderzulagen

Ein Bausparvertrag ist zuteilungsreif, wenn die Zuteilungsvoraussetzungen er-
füllt sind. In der Regel handelt es sich dabei um die Mindestvorgaben für das an-
gesparte Guthaben (das Mindestguthaben beträgt oft zwischen 40 \% und 50 \%
der Bausparsumme), für den Zeitraum, in der der Vertrag bereits besteht (Min-
destvertragsdauer), für den Zeitraum, in dem Einzahlungen geleistet wurden
(Mindestsparzeit) sowie für die vertraglich vereinbarte Bewertungszahl (Min-
destbewertungszahl). Ein zuteilungsreifer Bausparvertrag ermöglicht dem Bau-
sparer die Ausübung seines Rechtsanspruchs auf das Bauspardarlehen.

Aus Sicht des Bausparers ist eine Fortsetzung der Ansparphase dann besonders interessant, 
wenn die im Bausparvertrag fest vereinbarten Guthabenzinsen im Vergleich zum jeweiligen 
Marktzinsniveau hoch sind.

Die Finanzierung der individuellen Immobilieninvestitionen erfolgt durch die Zuteilung 
der im Bausparvertrag festgelegten Bausparsumme. Die Höhe des Bauspardarlehens ergibt sich 
als Differenz zwischen der Bausparsumme und dem Bausparguthaben zu Beginn der Darlehensphase. 
Die Gewährung des Darlehens ist typischerweise an die Stellung von Kreditsicherheiten gebunden.


Wie kann der Erwerb von Immobilien zur privaten Altersvorsorge genutzt werden.

Der Kauf von Immobilien während des Erwerbslebens kann aus Vorsorgeprodukt
zur Sicherung des Lebensstandards im Alter interpretiert werden. Grundsätzlich
sind hier die Alternativen des selbstgenutzten Wohnraums, der Vermietung von
Immobilien sowie der Erwerb von Immobilienfonds zu unterscheiden.



Diskutieren Sie die Eignung selbstgenutzten Wohneigentums zur Altersvorsorge.

Die Beurteilung der Eignung selbstgenutzten Wohneigentums zur Altersvorsorge ist komplex und einzelfallabhängig. 
Grundsätzlich handelt es sich dabei um
eine sehr langfristig wirkende Entscheidung mit weitgehend fixierten Rahmenbe-
dingungen ohne Diversifikationseffekte. Darüber hinaus erfolgt die Beurteilung der
eigenen Wohnung oder des eigenen Hauses in vielen Fällen nicht nur finanzrational
sondern unter zahlreichen emotionalen Gesichtspunkten.

Selbstgenutztes Wohneigentum als Anlageprodukt betrachten zu wollen, erfordert
eine möglichst vollständige Erfassung aller mit dem Erwerb, der Nutzung und der
Veräußerung anfallenden Zahlungen. Neben dem Kaufpreis der Immobilie und
den Erwerbsnebenkosten (Grunderwerbssteuer, Makler- und Notargebühren)
sind auch alle Finanzierungskosten (Darlehenszinsen und Gebühren) zu berück-
sichtigen. Während der Nutzung anfallende Instandhaltungs- und Renovie-
rungskosten sind auch zu erfassen. Der Summe dieser auf den Veräußerungs-
zeitpunkt aufgezinsten (Opportunitätszins) Auszahlungen ist dann der realisierte
Verkaufserlös abzüglich ggf. anfallender Nebenkosten gegenüber zu stellen.


Was ist eine Umkehrhypothek und wo liegen die wesentlichen Unterschiede zur Leibrente gem. §§ 759-761 BGB?

Ein möglicher Weg für die Senioren, den Verbleib in der eigenen Immobilie sicher-
zustellen besteht im Abschluss eines Kreditvertrages mit der Immobilie als Sicher-
heit (sog. Umkehrhypothek). Das Darlehen kann dann als Einmalbeitrag oder als
lebenslange Rente ausbezahlt werden, der Kreditnehmer bleibt Eigentümer der
Immobilie und kann diese weiterhin bewohnen. Zur Sicherstellung des Darlehens
wird eine Grundschuld bestellt, Zinsen und Tilgung werden typischerweise gestun-
det. Aus diesem Grund steigt die Schuldenlast im Zeitablauf solange bis sie nach
dem Tod des Darlehensnehmers oder bei dessen Umzug in eine Pflegeeinrichtung
von den Erben getilgt oder durch die Veräußerung der Immobilie ausgeglichen
wird.


Als weitere Gestaltungsmöglichkeit für Senioren in der eigenen Immobilie bleiben
zu können und zugleich eine Altersrente zu beziehen, entwickelt sich die Leibrente
(§§ 759-761 BGB). Dabei verkaufen die Eigentümer ihr Haus oder ihre Wohnung
und erhalten im Gegenzug ein lebenslanges Wohnrecht sowie eine monatliche
Leibrente.
Das mietfreie Wohnrecht und die monatliche Leibrentenzahlung werden
notariell beurkundet und im Grundbuch eingetragen.


Was verstehen wir unter dem Versicherungstechnischen Risiko? Wel-
che Komponenten können theoretisch unterschieden werden? Erläutern
Sie die Entstehungsgründe und grundlegende Messkonzepte des versi-
cherungstechnischen Risikos.

Wesensmerkmale und Entstehungsgründe: Das versicherungstechnische Risiko bezeichnet in erster Näherung die Gefahr, 
dass die Auszahlungen für Versicherungsfälle die korrespondierenden Prämieneinzahlungen überschreiten. 
Diese unter betriebswirtschaftlichen Aspekten zweckmäßige Wesensbeschreibung des
versicherungstechnischen Risikos betont neben der Unsicherheitsdimension
eine Finalitätsdimension.

Die Unsicherheit betrifft den Eintritt, den Zeitpunkt
und die Höhe künftiger Entschädigungszahlungen. Das betriebswirtschaftlich rele-
vante Risiko entsteht aufgrund des Erfordernisses, dieser zufallsabhängigen Scha-
denentwicklung a priori ein deterministisches Leistungsäquivalent gegenüber zu
stellen. Zum Zeitpunkt des Vertragsabschlusses und damit der Prämienvereinba-
rung ist die Hauptkostenkomponente des Vertrages der Versicherungsunterneh-
mung unbekannt.


Komponenten: Als theoretisch isolierbare Komponenten des versicherungstech-
nischen Risikos können Irrtums- und Zufallsrisiko unterschieden werden. Das Irr-
tumsrisiko resultiert aus der Unvollständigkeit der Informationen über die wahre
Zufallsgesetzmäßigkeit der Schäden und zerfällt in die Bestandteile des Diagnose-
und des Prognoserisikos. Das Diagnoserisiko besteht in der Gefahr, die in der
Vergangenheit gültige Zufallsgesetzmäßigkeit der versicherungstechnischen Ein-
heit nicht richtig zu identifizieren. Mögliche Ursachen hierfür liegen in einer fehler-
haften Modellauswahl und -spezifizierung sowie in der Unvollständigkeit der ver-
wendeten Daten (statistische Inferenz).


Das Prognoserisiko (statistische Prognose) resultiert aus der ex ante prinzipiell
nicht zu bestätigenden Hypothese über die Stabilität bzw. die konkrete Entwick-
lung der inferierten Gesetzmäßigkeit. Selbst bei angenommener fehlerfreier Diag-
nose besteht die Unsicherheit, ob die für die Vergangenheit festgestellte Schaden-
gesetzmäßigkeit auch in der Zukunft Gültigkeit besitzt.


Neben dem Irrtumsrisiko beschreibt das Zufallsrisiko die zweite Komponente des versicherungstechnischen Risikos, 
da selbst bei vollständiger Kenntnis der wahren Schadengesetzmäßigkeit die tatsächliche Rea-
lisation a priori unbekannt bleibt. Es verbleibt aufgrund der Natur des Zufalls stets
eine positive (Rest-)-Wahrscheinlichkeit dafür, dass die tatsächlich zu entrichten-
den Auszahlungen für Versicherungsleistungen nicht aus dem Gesamtbetrag der
zur Risikodeckung regelmäßig vorhandenen Vermögenswerte finanziert werden
können.



Skizzieren und systematisieren Sie die Finanzströme der Versicherungs-
unternehmung. 



Erläutern Sie den Zusammenhang zur Ermittlung des
Kapitalbedarfs, der Bildung von Kapital und der Investitionsentschei-
dung. Grenzen sie die Produktionsstufen der Privatversicherung ab und nen-
nen Sie zugehörige wesentliche versicherungstechnische Instrumente.

INSERT IMAGE 74


Die Identifikation und Bewertung vom Versicherungsnehmer auf das Versiche-
rungsunternehmen übertragenen Risikos ist Kern des versicherungstechnischen
Risikotransfers.


Unterscheiden Sie die zulässigen Rechtsformen der Versicherungsun-
ternehmen und begründen Sie diese Beschränkung.

Die Aufgabe des Aktuars besteht dabei in der Kalkulation der individuell risikoadä-
quaten Versicherungsprämie (Bruttorisikoprämie), die sich grundsätzlich aus
dem erwarteten Schaden und einem Sicherheitszuschlag (z.B. Kapitalkostenzu-
schlag) zusammensetzt. Der am Markt erzielbare Preis (Marktprämie) für den Ri-
sikotransfer muss nicht identisch mit der aktuariell kalkulierten Prämie sein. Wei-
tere Instrumente des Risikotransfers sind z.B. die Risikoselektion durch An-
nahme/Ablehnung von Risiken, die Tarifbildung sowie die individuelle Vertragsge-
staltung (materialer Deckungsumfang, Deckungsgrenzen und Selbstbeteiligun-
gen).


Kapitalreserven werden zunächst benötigt um in der Zukunft die erwarteten Schäden zu
decken (Rückstellungen). Um auch die zufällig über den Erwartungswert hinaus
eingetretenen Schäden ausgleichen zu können, bedarf es weiterer Kapitalreserven
(Sicherheitskapital).

Der Sicherheitskapitalbedarf reflektiert dabei die Gesamtrisikoposition des Versi-
cherers nach Risikotransformation, d.h. unter Berücksichtigung kollektiver und
zeitlicher Ausgleichseffekte. Zu den Instrumenten der Risikotransformation
zählen neben der Kapitalbildung auch die Vermögensanlagepolitik, die Rückversi-
cherungspolitik und die planmäßige Bestandsorganisation.



Skizieren Sie mögliche Kapitalbildungsmöglichkeiten von Versicherungen

Kapitalanlage
Finanzielle Mittel, die im Unternehmen verbleiben, werden an den Finanz- und Ka-
pitalmärkten investiert (versicherungstechnische Kuppelproduktion)

Mit der der
Solvency II-Umsetzung
wurde die ehemalige Anlageverordnung
(AnlV) mit Wirkung zum 1. Januar 2016 aufgehoben. An die Stelle der bisherigen
Anlagegrundsätze für das gebundene Vermögen nach § 54 Abs. 1 bis 3 VAG a. F.
sind die neuen Anlagegrundsätze des § 124 Abs. 1 VAG getreten.


Rückversicherung und Zahlungsströme


NOT SO IMPORTANT

Ermittlung des Rückversicherungssaldos (aus Sicht des Erstversicherers):
\( R_S \)
Rückversicherungssaldos
\( R_p \)
Rückversicherungsprämie (Aufwand)
\( S_R \)
Anteil des Rückversicherers an der Veränderung der Rückstellung für Bei-
tragsüberträge
(Erhöhung / Verminderung des RV-Anteils = Ertrag / Aufwand)
\( X_R \)
Anteil des Rückversicherers an den Schadenzahlungen (Ertrag)
\( X_C \)
Anteil des Rückversicherers an der Veränderung der Schadenrückstellung
(Erhöhung / Verminderung des RV-Anteils = Ertrag / Aufwand)
\( S_C \)
Anteil des Rückversicherers an den Betriebskosten (Ertrag)

\( R_S = - R_p \pm S_R + X_R \pm X_C + S_C \)

Risikokompensation

Die Risikokompensation bezeichnet die Konkretisierung des Schutzversprechens
im Versicherungsfall. Bei der Kompensation sind neben der reinen Zahlung des versicherten Schadens durch
das Versicherungsunternehmen zusätzlich Regulierungskosten (Prüfung des Schadens im Unternehmen bzw. 
durch externe Gutachter etc.) zu beachten. Sie setzt sich
aus den Bereichen Schadenerfassung, -bearbeitung und -regulierung zusammen.
Früher wurde hier typischerweise die substanziell nicht beeinflussbare Zahlungspflicht gegenüber dem Versicherungs-
nehmer betont. Heute erfolgt zunehmend die Hinwendung zu einem aktiven Schadenmanagement.



Unterscheiden Sie die zulässigen Rechtsformen der Versicherungsun-
ternehmen und begründen Sie diese Beschränkung.


§ 8 (2) des Versicherungsaufsichtsgesetzes (VAG) bestimmt, dass die aufsichts-
behördliche Erlaubnis zum Geschäftsbetrieb eines Versicherungsunternehmens
nur dann erteilt werden darf, wenn es sich dabei um eine 

- Versicherungsaktiengesellschaft (VersAG) 
- Europäischen Gesellschaft (Societas Europeae, SE), 
- Versicherungsverein auf Gegenseitigkeit (VVaG)
- öffentlich-rechtliches Versicherungsunternehmen (ÖRVU) 
handelt.

Als juristischen Personen unterscheiden sie sich insb. hinsichtlich ihrer Träger-
schaft und Organbezeichnung sowie hinsichtlich der jeweiligen Rechtsgrund-
lagen, mit entsprechenden Auswirkungen insb. auf ihre Möglichkeiten zur exter-
nen Kapitalbeschaffung und passiven Konzernfähigkeit.

Dem aufsichtsrechtlichen Schutzgedanken folgend ist es nicht zulässig, Versiche-
rungsunternehmen beispielsweise als Gesellschaft mit beschränkter Haftung
(GmbH) oder als Personengesellschaft wie einer Offenen Handelsgesellschaft
(OHG) oder einer Kommanditgesellschaft (KG) zu betreiben.



Diskutieren Sie die Angleichung der Unternehmensformen im Wettbe-
werb

Unabhängig von ihrer Rechtsform ist das Management der Versicherungsunterneh-
men im Wettbewerb gleichgestellt. Dies betrifft zunächst vor allem die rechtsform-
unspezifische Versicherungsaufsicht (Zulassung, Solvabilität, Kapitalanlage) und
ebenso das jeweils nicht-spezifische Versicherungsvertrags-, Vermittler-, Wettbe-
werbs- und Unternehmenssteuerrecht.Die Zusammenarbeit in Versicherungsver-
bänden ist ebenso rechtsformübergreifend, wie die Verwendung identische versi-
cherungstechnischer und betriebstechnischer Verfahren.
Diskutieren Sie die Angleichung der Unternehmensformen im Wettbe-
werb


Unterscheiden Sie die Organe der Versicherungsunternehmen und erläutern Sie wesentliche Aufgaben.
Erläutern Sie den „fit and proper“-Grundsatz

Gesetzliche Organe der Versicherungsaktiengesellschaft (VersAG) sind: Vorstand,
Aufsichtsrat und Hauptversammlung. Die Vertretung der Arbeitnehmerinteressen
erfolgt durch die von der Belegschaft gewählten Arbeitnehmervertreter im mitbe-
stimmten Aufsichtsrat. Die gesetzlichen Organe des Versicherungsvereins auf Ge-
genseitigkeit (VVaG) sind analog strukturiert und haben grundsätzlich identische
Aufgaben. Anstelle der Hauptversammlung, auf der die Aktionäre ihre Interessen
wahrnehmen, tritt hier das Oberste Organ als Versammlung der Mitglieder oder
als Mitgliedervertreterversammlung. Bei öffentlich-rechtlichen Versicherungsun-
ternehmen (ÖRVU) unterscheiden wir lediglich Vorstand (Geschäftsführungsorgan)
und Verwaltungsrat.

Der Vorstand besteht aus mindestens 2 Mitgliedern, er leitet das Unternehmen in
Vorstand verantwortet die Aufstellung des Jahresabschlusses und die Einberufung
der Hauptversammlung, resp. der Mitglieder- bzw. Mitgliedervertreterversamm-
lung. Er hat dafür Sorge zu tragen, dass die Geschäftsorganisation regelmäßig
intern überprüft wird. Der Vorstand ist zur regelmäßigen Berichterstattung gegen-
über dem jeweiligen Aufsichtsorgan verpflichtet. Die Bestellung als Vorstandsmit-
glied ist auf max. 5 Jahre beschränkt. Die Bafin prüft für jedes Mitglied des Vor-
standes, dessen Zuverlässigkeit und fachliche Eignung („fit and proper“).
eigener Verantwortung und er vertritt die Gesellschaft im Außenverhältnis.


Der Aufsichtsrat, resp. der Verwaltungsrat bestellt den Vorstand und ist ver-
pflichtet, die Geschäfts-führung des Vorstands zu überwachen. Zu diesem Zweck
erhält der Aufsichtsrat regelmäßig die erforderlichen Berichte, z.B. ORSA-Bericht.
Die Satzung bestimmt darüber hinaus, Rechtsgeschäfte von grundlegender Bedeu-
tung, die der Zustimmung des Aufsichtsrates bedürfen. Der Aufsichtsrat beauftragt
die Abschlussprüfer und ist für die Prüfung und Feststellung des Jahresabschlusses
verantwortlich. Aufsichtsratsmitglieder müssen zuverlässig sein und die zur Wahr-
nehmung der Kontrollfunktion sowie zur Beurteilung und Überwachung der Ge-
schäfte, die das Unternehmen betreibt, erforderlichen Fähigkeiten und Erfahrun-
gen nachweisen (Sachkundenachweis).
